% ============================================================
%  intro.tex — 引言:motivation → gap → 贡献点(编号,可证伪)
%  纪律:主张回指来源(见 traceability-worksheet.md 主张台账);比较级须有支撑。
% ============================================================

\section{Introduction}
\label{sec:intro}

% [动机] 从领域真实痛点切入,不空泛。每条事实性主张 \cite{}。
Dexterous manipulation is central to embodied intelligence because \ldots~\cite{<key>}.
However, current approaches struggle with \ldots~\cite{<key>}.

% [gap] 明确指出现有工作的具体缺口(与 ① GapScout 的差异化陈述一致)。
Prior work~\cite{<key>,<key>} addresses \ldots, yet \textbf{remains limited in}
\ldots. This gap matters because \ldots

% [本文思路] 一句话点出核心 insight / 方法直觉。
In this paper, we argue that \ldots and propose \textbf{MethodName}.

% [贡献点] 编号、可证伪、每条对应后文实验。禁「首个/最优」等无支撑最高级。
Our contributions are:
\begin{itemize}
  \item \textbf{C1:} A method that \ldots (validated in Sec.~\ref{sec:exp}, Table~\ref{tab:main}).
  \item \textbf{C2:} An ablation showing \ldots (Sec.~\ref{sec:ablation}).
  \item \textbf{C3:} Real-robot results demonstrating \ldots (Sec.~\ref{sec:real}, 采集条件见附录).
\end{itemize}

% [路线图] 可选:简述后文结构。
